\documentclass[11pt]{article}

\usepackage[T1]{fontenc}
\usepackage{lmodern}
\usepackage[margin=1in]{geometry}
\usepackage{amsmath,amssymb,amsthm}
\usepackage{booktabs}
\usepackage[hidelinks]{hyperref}

\newtheorem{theorem}{Theorem}
\newtheorem{lemma}[theorem]{Lemma}
\newtheorem{corollary}[theorem]{Corollary}
\theoremstyle{remark}
\newtheorem{remark}[theorem]{Remark}

\newcommand{\F}{\mathbb{F}}

\title{New lower bounds for mixed binary/ternary covering codes}
\author{Noah Parsons\\ \small Independent researcher\thanks{%
\texttt{parsons.m.noah@gmail.com}, ORCID \href{https://orcid.org/0009-0000-7224-6040}{0009-0000-7224-6040}.}}
\date{October 2026}

\begin{document}
\maketitle

\begin{abstract}
Let $K(t,b,1)$ be the minimum size of a code in $\F_3^t\times\F_2^b$ with covering radius $1$.
We prove $K(2,5,1)\ge 34$, $K(1,7,1)\ge 45$ and $K(2,6,1)\ge 59$, improving the lower bounds
$32$, $42$ and $57$ in the tables of mixed binary/ternary covering codes. These are three of the
five smallest open entries for covering radius $1$. The proofs follow the codeword-distribution
method of Östergård and Blass: the space is divided into cells, a counting inequality restricts
the number of codewords in each cell, and every remaining distribution is shown infeasible by
integer programming. Each case is decided by two independent solvers, Gurobi and CPLEX, and
the steps that produce the cases are checked against the definitions by separately written code.
\end{abstract}

\section{Introduction}

A code $C\subseteq\F_3^t\times\F_2^b$ has \emph{covering radius $1$} if every word is within
Hamming distance $1$ of some codeword. The minimum size of such a code is denoted
$K(t,b,1)$~\cite{cohen-book}. Bounds for $t+b\le 13$ are collected in the tables of
K\'eri~\cite{keri}, which follow the tables of Östergård and Hämäläinen~\cite{oh} and of
Bertolo, Östergård and Weakley~\cite{bow}; the tables were last updated in 2011.

Ordered by the size $3^t2^b$ of the space, the five smallest open entries for covering radius $1$
are $K(2,5,1)$, $K(1,7,1)$, $K(3,4,1)$, $K(5,1,1)$ and $K(2,6,1)$ (spaces of $288$, $384$, $432$,
$486$ and $576$ words). This note improves the lower bound of three of them.

\begin{theorem}\label{thm:main}
$K(2,5,1)\ge 34$, $K(1,7,1)\ge 45$ and $K(2,6,1)\ge 59$.
\end{theorem}

\begin{table}[h]
\centering
\begin{tabular}{lccc}
\toprule
Entry & Previous bounds~\cite{keri} & Sources & New lower bound \\
\midrule
$K(2,5,1)$ & $32$--$36$ & \cite{lqs}, \cite{hr} & $34$ \\
$K(1,7,1)$ & $42$--$48$ & \cite{lqs}, $3K(0,7,1)$ & $45$ \\
$K(2,6,1)$ & $57$--$64$ & \cite{lqs}, \cite{hr} & $59$ \\
\bottomrule
\end{tabular}
\caption{The three entries. All three previous lower bounds are due to Lang, Quistorff and
Schneider~\cite{lqs}. The upper bounds are due to Hämäläinen and Rankinen~\cite{hr}, and for
$K(1,7,1)$ to the direct product of $\F_3$ with the binary Hamming code of length $7$.}
\label{tab:entries}
\end{table}

We are not aware of improvements to these entries after 2011. Since a covering code with fewer
than $M$ words can be enlarged to exactly $M$ words, a lower bound $K\ge M+1$ follows from the
nonexistence of covering codes with exactly $M'$ words for each $M'$ from the previous lower bound
up to $M$.

\section{Method}

\subsection{Cells and the counting inequality}

Fix a set $T$ of coordinates. For each assignment $x$ of values to the coordinates in $T$, the
\emph{cell} $S_x$ is the set of words that agree with $x$ on $T$. It has $2^f$ words when the
free coordinates (those outside $T$) are $f$ binary coordinates, as in all cases below. Two cells
are \emph{adjacent} if their assignments differ in exactly one coordinate. For a code $C$ write
$M_x=|C\cap S_x|$.

\begin{lemma}\label{lem:count}
If $C$ has covering radius $1$ and the free coordinates are $f$ binary coordinates, then for every
cell $S_x$
\[
(1+f)\,M_x+\sum_{y\ \text{adjacent to}\ x} M_y\ \ge\ 2^{f}.
\]
\end{lemma}

\begin{proof}
The balls of radius $1$ around the codewords cover the $2^f$ words of $S_x$. A codeword in $S_x$
covers itself and the $f$ words of $S_x$ obtained by changing one free coordinate. A codeword in
an adjacent cell differs from $S_x$ in exactly one fixed coordinate, so its ball meets $S_x$ in
exactly one word. A codeword in any other cell differs from every word of $S_x$ in at least two
coordinates, so its ball misses $S_x$.
\end{proof}

Lemma~\ref{lem:count} is the inequality of Östergård and Blass~\cite{ob}, used there for
$K(9,1)=62$, written for mixed alphabets. We use two kinds of cells.

\paragraph{Two ternary coordinates.} For $K(2,b,1)$ we fix both ternary coordinates. The nine
cells form a $3\times 3$ grid, each cell has $f=b$ free coordinates, and a cell is adjacent to the
four cells in its row and column. Writing $R$ and $C$ for the totals of the row and column of a
cell, the inequality becomes
\begin{equation}\label{eq:grid}
(b-1)\,M_x+R+C\ \ge\ 2^b ,
\end{equation}
since $M_x$ is counted once in $R$ and once in $C$.

\paragraph{One ternary coordinate.} For $K(1,7,1)$ we fix the ternary coordinate and then the
binary coordinates $v_1,v_2,\dots$ one at a time, as in~\cite{ob}. At level $m$ there are $3\cdot
2^m$ cells with $f=7-m$; a cell $(a,x)$ is adjacent to the two cells $(a',x)$, $a'\neq a$, and to
the $m$ cells $(a,x+e_i)$. Each distribution at level $m+1$ refines one at level $m$, and the
inequality is imposed at every level.

\subsection{Symmetry and exact cases}

Permuting the symbols of a ternary coordinate, permuting coordinates of equal alphabet size, and
complementing binary coordinates are isometries, so equivalent distributions need only be
considered once. For the grid these maps act as permutations of rows, permutations of columns and
transposition, a group of order $72$. For the levels of $K(1,7,1)$ at level $2$ they act as the
$3!$ permutations of the ternary symbols and the $2^2\cdot 2!$ signed permutations of $v_1,v_2$,
a group of order $48$.

Each remaining distribution becomes an integer program with one binary variable per word, one
covering constraint per word (the sum of the variables in its ball is at least $1$) and one
constraint per cell fixing the number of codewords in it. A covering code with $M$ words exists
if and only if at least one of these programs is feasible.

\section{Computation}

All computations were run single-threaded on one core of an Intel Core Ultra~7 255HX, with
Gurobi~13.0 (gurobipy, size-limited licence) and IBM CPLEX~22.2 (Community Edition). The largest
models, for $K(2,6,1)$, have $576$ variables and $585$ constraints, within both licences.

\begin{table}[h]
\centering
\small
\begin{tabular}{lrrrccl}
\toprule
Entry & $M$ & distrib. & cases & Gurobi & CPLEX & further solvers \\
\midrule
$K(2,5,1)$ & 32 & 306     & 10  & all inf. & all inf. & SCIP, HiGHS: all inf. \\
           & 33 & 1\,194  & 32  & all inf. & all inf. & SCIP: all; HiGHS: 31 of 32 \\
\midrule
$K(1,7,1)$ & 42 & 596     & 24  & all inf. & all inf. & SCIP: all inf. \\
           & 43 & 3\,936  & 109 & all inf. & all inf. & SCIP: all inf. \\
           & 44 & 18\,108 & 489 & all inf. & all inf. & -- \\
\midrule
$K(2,6,1)$ & 57 & 372     & 14  & all inf. & all inf. & -- \\
           & 58 & 2\,637  & 64  & all inf. & all inf. & -- \\
\bottomrule
\end{tabular}
\caption{The computation. ``Distrib.'' counts the distributions over the cells (for
$K(1,7,1)$, at level $2$) that satisfy Lemma~\ref{lem:count}; ``cases'' counts them up to
symmetry, the integer programs solved; ``inf.'' means infeasible. One HiGHS case for $K(2,5,1)$ was still running after 30 minutes and was not
completed; SCIP was stopped after part of the $K(1,7,1)$ cases at $M=44$ for lack of speed.}
\label{tab:computation}
\end{table}

Table~\ref{tab:computation} lists the cases. Every case of every entry was found infeasible by
both Gurobi and CPLEX. The longest single case took $54$ seconds with Gurobi and $10$ minutes with
CPLEX (both for $K(2,6,1)$ at $M=58$); the total for all cases was under $30$ minutes with Gurobi
and under one hour with CPLEX. SCIP~10.0 and HiGHS (through SciPy~1.16.3) were used as further
checks where their speed allowed.

\section{Verification}

A false ``no code exists'' could come from a wrong counting inequality, an incomplete list of
distributions, a faulty symmetry reduction, a wrong integer program, or a solver error. Each was
checked separately.
\begin{enumerate}
\item \textbf{The inequality from the definition.} For every cell $S_x$ and every word $u$, the
  number of words of $S_x$ whose ball contains $u$ was computed from Hamming distances alone. In
  every case it depends only on the cell of $u$ and equals the coefficient in
  Lemma~\ref{lem:count}: $1+f$ for $u\in S_x$, $1$ for adjacent cells and $0$ otherwise. This was
  done for the grid with $b=5,6$ and for levels $0$, $1$, $2$ of $K(1,7,1)$.
\item \textbf{Completeness of the distributions.} For $K(2,5,1)$, a separate vectorised
  enumeration of all $76\,904\,685$ ways to place $32$ words in nine cells ($95\,548\,245$ for
  $33$), using the coefficients computed in step~1, gives exactly the $306$ ($1\,194$)
  distributions of Table~\ref{tab:computation}. For $K(1,7,1)$, a separate brute-force enumeration
  through levels $0$, $1$, $2$ gives exactly $596$, $3\,936$ and $18\,108$. For $K(2,6,1)$, where
  a full enumeration is too slow, the case lists come from a fast enumerator based on
  \eqref{eq:grid}; on seven smaller instances ($b=3,4,5$) its output is identical to that of the
  direct brute-force enumerator.
\item \textbf{Symmetry reduction.} With the group elements written out by separate code, every
  distribution in step~2 is equivalent to one of the cases solved, and every case solved occurs.
\item \textbf{The integer program.} A real covering code satisfies the program for its own cell
  distribution, and the same code with one word removed does not: a $36$-word code for
  $K(2,5,1)$, the $48$-word record code for $K(1,7,1)$, and a $74$-word code for $K(2,6,1)$
  (all confirmed by an independent brute-force checker).
\item \textbf{Known values.} The same pipelines reproduce $K(2,3,1)=12$, $K(2,4,1)=20$,
  $K(1,5,1)=16$ and $K(1,6,1)=24$: each reports no code one below the known value and finds a code
  at the known value, confirmed by the checker. Gurobi and CPLEX each pass these controls.
\item \textbf{Solvers.} Every case was decided by two solvers from different developers. The models
  for $K(2,5,1)$ and $K(2,6,1)$ were built by separate code for each solver; for $K(1,7,1)$ both
  solvers received models built by the same code, which step~4 tests.
\end{enumerate}

\begin{remark}
All the solvers use floating-point arithmetic, so each verdict is evidence rather than an exact
certificate. The independent agreement of two commercial-grade solvers on every case, together
with the checks above, is the basis of Theorem~\ref{thm:main}.
\end{remark}

\section{Remarks}

\begin{enumerate}
\item For $K(1,7,1)$ we showed separately~\cite{companion} that any covering code with at most
  $47$ words has trivial automorphism group. With Theorem~\ref{thm:main}, a code improving the
  upper bound $48$ would have $45$, $46$ or $47$ words and no symmetry.
\item The method continues to the next sizes: $K(2,5,1)$ at $M=34$ has $93$ cases up to symmetry,
  $K(1,7,1)$ at $M=45$ has $1\,848$, and $K(2,6,1)$ at $M=59$ has $221$. These were not
  completed.
\item The entries $K(3,4,1)$ and $K(5,1,1)$ have three or more ternary coordinates, which the
  programs used here do not handle; they were not attempted.
\end{enumerate}

\paragraph{Code and data.} All programs and logs are in the public repository
\begin{center}\url{https://github.com/Noah-Parsons/covering-codes}\end{center}
\begin{itemize}
\item grid method: \texttt{src/split\_t2.py};
\item its rechecks: \texttt{src/recheck\_t2\_gurobi.py},
\item[] \hphantom{its rechecks: }\texttt{src/recheck\_t2\_cplex.py}, \texttt{src/recheck\_t2\_highs.py};
\item its checks: \texttt{src/verify\_t2\_distributions.py};
\item level method: \texttt{src/exact\_dist.py}, \texttt{src/leaves\_t1.py};
  checks \texttt{src/verify\_t1\_leaves.py};
\item brute-force code checker: \texttt{src/checker.py}; logs in \texttt{results/}.
\end{itemize}

\begin{thebibliography}{9}

\bibitem{bow} R.~Bertolo, P.~R.~J. Östergård and W.~D. Weakley, An updated table of
binary/ternary mixed covering codes, \emph{J. Combin. Des.} 12 (2004), 157--176.

\bibitem{cohen-book} G.~Cohen, I.~Honkala, S.~Litsyn and A.~Lobstein, \emph{Covering Codes},
North-Holland, Amsterdam, 1997.

\bibitem{hr} H.~Hämäläinen and S.~Rankinen, Upper bounds for football pool problems and mixed
covering codes, \emph{J. Combin. Theory Ser.~A} 56 (1991), 84--95.

\bibitem{keri} G.~Kéri, Tables for bounds on covering codes,
\url{https://old.sztaki.hu/~keri/codes/index.htm} (accessed 16 October 2026).

\bibitem{lqs} W.~Lang, J.~Quistorff and E.~Schneider, New results on integer programming for
codes, \emph{Congr. Numer.} 188 (2007), 97--107.

\bibitem{ob} P.~R.~J. Östergård and U.~Blass, On the size of optimal binary codes of length $9$
and covering radius $1$, \emph{IEEE Trans. Inform. Theory} 47 (2001), 2556--2557.

\bibitem{oh} P.~R.~J. Östergård and H.~O. Hämäläinen, A new table of binary/ternary mixed
covering codes, \emph{Des. Codes Cryptogr.} 11 (1997), 151--178.

\bibitem{companion} N.~Parsons, Improving the covering code bound $K(1,7,1)\le 48$ requires an
asymmetric code, manuscript, 2026.

\end{thebibliography}
\end{document}
